\begin{lemma}[A classically differentiable function lies in $W^{1,q}_{\mathrm{loc}}$ with its classical gradient]
Let $n\ge1$ be an integer and let $q\ge1$ be real.  Let
$\Omega\subseteq\R^n$ be open, let $u\colon\R^n\to\R$ and let
$g\colon\R^n\to\R^n$.  Assume
\begin{enumerate}
\item $u$ is differentiable at every point $x\in\Omega$, and for every such $x$
  and every $i\in\{1,\dots,n\}$,
  \[
  D_iu(x)\;=\;g^i(x),
  \]
  where $D_iu(x)=Du(x)[e_i]$ is the derivative of $u$ at $x$ in the direction of
  the $i$-th standard basis vector;
\item $g$ is continuous on $\Omega$.
\end{enumerate}
Then $u\in W^{1,q}_{\mathrm{loc}}(\Omega)$ with weak gradient $g$
(\noderef{MemW1pLocOn}).

\noindent\textbf{On the hypothesis shape.}  No regularity of $u$ or $g$ off
$\Omega$ is assumed, and none is needed: both clauses of
\noderef{MemW1pLocOn} are statements about $\Omega$ and about compact subsets of
$\Omega$.  Clause (1) is exactly the first clause of the tablet's
$C^{1,\beta}$ predicate, and clause (2) is the continuity that the tablet's
continuous-representative lemma supplies for the field it produces, so the $C^1$ representative of \S4 and \S5
feeds this lemma without reshaping.  Clause (2) is not implied by clause (1) --
a differentiable function need not have continuous derivative -- and it is what
makes the three integrands below continuous.  Continuity of $u$ on $\Omega$ is
\emph{not} assumed separately: it follows from clause (1), since a function
differentiable at a point is continuous there.
\end{lemma}
\begin{proof}
Write $e_i$ for the $i$-th standard basis vector of $\R^n$ and
$D_ih(x)=Dh(x)[e_i]$ throughout.  By \noderef{MemW1pLocOn} there are two things
to prove: the weak-gradient identity of \noderef{HasWeakGradientOn} for the pair
$(u,g)$ on $\Omega$, and, for every compact $K\subseteq\Omega$, that $u$ and $g$
lie in $L^q(K)$.

\emph{Step 0: $u$ is continuous on $\Omega$, and $g^i$ is continuous on
$\Omega$ for each $i$.}  By clause (1) $u$ is differentiable at each $x\in\Omega$,
hence continuous at each $x\in\Omega$, so $u$ is continuous on $\Omega$.  By
clause (2) $g$ is continuous on $\Omega$, and each coordinate map
$h\mapsto h^i$ on $\R^n$ is continuous (it is linear on a finite-dimensional
space), so $g^i=(\,\cdot\,)^i\circ g$ is continuous on $\Omega$ for each $i$.
Finally, by clause (1), the function $x\mapsto D_iu(x)$ agrees with $g^i$ at
every point of $\Omega$, so it too is continuous on $\Omega$.

\emph{Step 1: the three integrands are continuous on all of $\R^n$ and
compactly supported.}  Fix $\varphi\in C^\infty(\R^n)$ with compact support and
$\operatorname{supp}\varphi\subseteq\Omega$, and fix an index $i$; by
\noderef{HasWeakGradientOn} these are exactly the test data, and
$\operatorname{supp}$ is the closed support throughout.  Put
\[
P_1(x)=D_iu(x)\,\varphi(x),\qquad
P_2(x)=u(x)\,D_i\varphi(x),\qquad
P_3(x)=u(x)\,\varphi(x).
\]
Let $U=\R^n\setminus\operatorname{supp}\varphi$, which is open because
$\operatorname{supp}\varphi$ is closed.  Since
$\operatorname{supp}\varphi\subseteq\Omega$ we have $\Omega\cup U=\R^n$, and
both sets are open.

On $\Omega$: $u$, $x\mapsto D_iu(x)$, $\varphi$ and $D_i\varphi$ are all
continuous --- the first two by Step 0, the last two because
$\varphi\in C^\infty(\R^n)$ --- so $P_1,P_2,P_3$ are continuous on $\Omega$ as
products of continuous functions.

On $U$: $\varphi$ vanishes identically on the open set $U$, hence so does its
derivative, so $D_i\varphi$ vanishes on $U$ as well.  Therefore
$P_1\equiv P_2\equiv P_3\equiv0$ on $U$, and each is continuous there.  (In
particular $P_1$ is unambiguous on $U$ whatever value is assigned to $D_iu$ at a
point where $u$ is not differentiable, because the factor $\varphi$ is $0$
there; the same remark applies to the integrals below.)

Continuity is a local property and $\Omega\cup U=\R^n$, so $P_1,P_2,P_3$ are
continuous on $\R^n$.  Moreover
$\{x:P_k(x)\ne0\}\subseteq\{x:\varphi(x)\ne0\}\cup\{x:D_i\varphi(x)\ne0\}
\subseteq\operatorname{supp}\varphi$ for each $k$ --- for $k=1,3$ because the
factor $\varphi$ must be nonzero, for $k=2$ because the factor $D_i\varphi$
must be, and $D_i\varphi$ vanishes off $\operatorname{supp}\varphi$ by the
previous paragraph --- and $\operatorname{supp}\varphi$ is closed, so
$\operatorname{supp}P_k\subseteq\operatorname{supp}\varphi$, which is compact.
Each $P_k$ is thus continuous with compact support, hence Lebesgue integrable on
$\R^n$.

\emph{Step 2: integration by parts.}  Apply the Preamble's integration by parts
in a direction with $f=u$, $g=\varphi$ and $v=e_i$, so that $D_v=D_i$.  Its
hypotheses hold:
\begin{itemize}
\item $f=u$ is differentiable at every point of $\operatorname{supp}\varphi$,
  because $\operatorname{supp}\varphi\subseteq\Omega$ and clause (1) gives
  differentiability at every point of $\Omega$;
\item $g=\varphi$ is differentiable at every point of $\operatorname{supp}u$,
  because $\varphi\in C^\infty(\R^n)$ is differentiable at every point of
  $\R^n$;
\item the three functions $x\mapsto D_iu(x)\varphi(x)$,
  $x\mapsto u(x)D_i\varphi(x)$ and $x\mapsto u(x)\varphi(x)$ --- that is,
  $P_1$, $P_2$ and $P_3$ --- are Lebesgue integrable on $\R^n$, by Step 1.
\end{itemize}
The conclusion is
\begin{equation}
\int_{\R^n}u\,D_i\varphi\,dx\;=\;-\int_{\R^n}D_iu\,\varphi\,dx .
\tag{1}
\end{equation}

\emph{Step 3: passing from $\R^n$ to $\Omega$.}  Both integrands in (1) vanish
off $\operatorname{supp}\varphi\subseteq\Omega$, by Step 1; so each integral
over $\R^n$ equals the corresponding integral over $\Omega$, and
\[
\int_\Omega u\,D_i\varphi\,dx\;=\;-\int_\Omega D_iu\,\varphi\,dx .
\]
On $\Omega$ clause (1) gives $D_iu=g^i$ pointwise, so the right-hand integrand
equals $g^i\varphi$ there and
\[
\int_\Omega u\,D_i\varphi\,dx\;=\;-\int_\Omega g^i\,\varphi\,dx .
\]
Since $\varphi$ and $i$ were an arbitrary test function and index, this is
precisely the defining identity of \noderef{HasWeakGradientOn} for $(u,g)$ on
$\Omega$.

\emph{Step 4: the local $L^q$ bounds.}  Let $K\subseteq\Omega$ be compact.  By
Step 0, $u$ and each $g^i$ are continuous on $\Omega$, hence on $K$; a
continuous real function on a compact set is bounded, so there is $M\ge0$ with
\[
|u(x)|\le M\quad\text{and}\quad |g(x)|\le M\qquad(x\in K),
\]
the second after taking $M$ larger than each of the $n$ coordinate bounds and
using $|g|\le\sum_i|g^i|$.  Also $K$ is compact, hence bounded, hence of finite
Lebesgue measure, $|K|<\infty$.  Finally $u$ and $g$ are continuous on $K$ and
$K$ is a Borel set, so their restrictions to $K$ are Borel measurable.  Hence
\[
\int_K|u|^q\,dx\le M^q|K|<\infty,\qquad
\int_K|g|^q\,dx\le M^q|K|<\infty,
\]
so $u\in L^q(K)$ and $g\in L^q(K)$, using $q\ge1$ so that the exponent is
admissible.  As $K$ was an arbitrary compact subset of $\Omega$, the second
clause of \noderef{MemW1pLocOn} holds.

Steps 3 and 4 together are the two clauses of \noderef{MemW1pLocOn}, so
$u\in W^{1,q}_{\mathrm{loc}}(\Omega)$ with weak gradient $g$.
\end{proof}
